\section{An exponential Padé estimate}\label{sec:pade}

Let \(F\) be a monic irreducible common factor of the reciprocal
moment polynomials \(R_{N-1},R_N\). Write \(h=\deg F\), \(H=dh\),
\(b=a+d\), \(M=bN\), \(m=N-1\), \(s=d^2/b\), and \(\alpha=d/b\).
We use the root bounds \(ai<|x|^{1/d}\le bi\) for \(R_i(x)=0\)
with \(i=m,N\), from Lemma~\ref{lem:gamma-annulus}. In particular,
\(H\log(aN)<\log|F(0)|<H\log M\). The Padé construction works on
\(z^d=-x\), where the root bound is linear in the moment order.
Its key feature is that an integral polynomial approximates the
same polynomial to high order after multiplication by each of the
\(d\) exponentials \(e^{-\omega z}\), \(\omega^d=1\).
The resulting Wronskian makes at least one of the approximants
nonzero above every root of \(F\). The comparison of an integer
resultant with its local estimates adapts the Laguerre construction
in Nguyen-Dang's preprint \cite[Proposition~3.9]{nguyen-dang}.
The following proof derives the rotated truncations, their
Wronskian and both local estimates for the binomial moments.

\begin{lemma}[Padé resultant inequality]\label{lem:resultant}
Assume \(N\ge200\), \(a,d\ge1\) and \(d<3a\). Set
\(r=(b-1)m\), and put
\[
 \begin{gathered}
 K_\pm=\tfrac12\pm\tfrac{\alpha}{14},\qquad
 \rho_0=b(1-\alpha K_+)=a+dK_-,\\
 g=\max\{b\log(b/\rho_0)-b+\rho_0,\,
          \rho_0-a-a\log(\rho_0/a)\}.
 \end{gathered}
\]
There is a monic irreducible factor \(G\) of \(F(-z^d)\) of degree
\(D=hk\), where \(1\le k\le d\), and \(K=|G(0)|\) satisfies
\(K^d=|F(0)|^k\). Write \(V=\log(M+1)\).
At each complex root \(z\) of \(G\) one may take the gain
\[
 L(z)=\max\left\{(m!e^{-mg-1})^{1/(m+1)},
                    \frac{am-\Re z}{d}\right\}.
\]
With \(\Lambda=(\prod_{G(z)=0}L(z))^{1/D}\), one has
\begin{equation}\label{eq:pade-norm}
 \log\Lambda<
 \sum_{p\mid K}\min\left\{\frac{\vp(K)}{D},\frac1{p-1}\right\}\log p
 +\frac{\omega(K)V}{N}.
\end{equation}
Here \(\omega(K)\) is the number of distinct prime divisors.
In addition, the first entry in \(L(z)\) can use \(g=0\)
when \(\Re z\le am\). For \(N\ge200\) the first entry is
greater than \(me^{-1.01-g}\). For \(d=1\) the additional gain
\(m\sqrt a\,e^{-1/2}/(1+a^{-1/2})\) is available.
\end{lemma}

\begin{proof}
Let \(B(z)=(am)!R_m(-z^d)/m!\) and let \(Q(z)\) be the Taylor
polynomial of \(e^z B(z)\) through degree \(r\). For \(0\le j\le bm\)
and \(v=bm-j\), direct expansion gives
\begin{equation}\label{eq:pade-coeff}
 [z^j]e^zB(z)=\frac{v!}{m!}
 [w^v](1+w)^{am}((1+w)^d-1)^m.
\end{equation}
Indeed the term indexed by \(i\) in \(R_m\) contributes
\((-1)^{m-i}\binom mi(am+di)!/(m!(am+di-v)!)\),
taken as zero if \(am+di<v\). Thus \(Q\) is integral, its top
coefficient is \(d^m\), and the coefficients of degrees \(r+1\)
through \(bm\) vanish.

For \(\omega^d=1\) put \(E_\omega(z)=e^{-\omega z}Q(\omega z)\).
Let \(T=e^{-z}Q-B\), which vanishes to order at least \(bm+1\).
Since \(B(\omega z)=B(z)\), we have \(E_\omega-B=T(\omega z)\).
Form the Wronskian of \(B\) and the \(d\) functions \(E_\omega\)
with derivative orders \(0,\ldots,d\). Factor \(e^{-\omega z}\)
out of each \(E_\omega\) column. The resulting polynomial
determinant has degree at most \(dm+dr=dbm\).
Its top coefficient comes from row zero of the \(B\) column
and the determinant \(((-\omega)^i)_{1\le i\le d,\,\omega^d=1}\),
which is a Vandermonde determinant times
\(\prod_{\omega^d=1}(-\omega)\). The top coefficient is nonzero.

To calculate the vanishing order, subtract the \(B\) column
from the original \(E_\omega\) columns. The invertible
change of columns taking the sums
\(d^{-1}\sum_{\omega^d=1}\omega^{-(bm+j)}T(\omega z)\),
for \(1\le j\le d\), retains one residue class of exponents in
each column. Their orders are at least \(bm+j\).
In each determinant term the derivatives subtract at most
\(\sum_{i=0}^d i\) from the total order, so the Wronskian vanishes
to order at least \(dbm\). Removing the exponential factors
preserves this order. Hence the polynomial determinant is a
nonzero multiple of \(z^{dbm}\). At any nonzero
root of \(B\) the zeroth row consequently has
\(Q(\omega z)\ne0\) for at least one \(\omega\).
Choose \(G\) to be the minimal polynomial of such a rotation
above a root of \(F\). Then \(\Res(G,Q)\) is a nonzero integer,
and every root of \(F\) has exactly \(k=D/h\) lifts in \(G\).
The norm identity follows by multiplying \(z^d=-x\) over them.

For \(\ell\ge1\), a beta integral gives the tail coefficient
\[
 [z^{bm+\ell}]e^z B(z)
 =\frac{(-1)^m}{m!(\ell-1)!}
       \int_0^1 t^{am}(1-t^d)^m(1-t)^{\ell-1}\,dt.
\]
With \(c_v=[t^v](1+t+\cdots+t^{d-1})^m\), the same coefficient is
\begin{equation}\label{eq:pade-beta}
 (-1)^m\binom{m+\ell-1}{m}
       \sum_v c_v \frac{(am+v)!}{(am+v+m+\ell)!}.
\end{equation}
Let \(p\) be prime and \(e=v_p(z)>0\) in an embedding into
\(\overline{\Q}_p\), normalized by \(v_p(p)=1\), and set
\(\mu=1/(p-1)\). The base-\(p\) digit sum \(s_p(n)\) satisfies
\(s_p(n)/(p-1)\le\log_p(n+1)\), since
\(p^{u/(p-1)}\le u+1\) for each digit \(u\), and the product
of the \(u+1\) counts integers whose digits do not exceed
those of \(n\). If \(e\le\mu\), the coefficient of \(z^{r-j}\) in
\eqref{eq:pade-coeff} contains the integer factor \((m+j)!/m!\).
By Legendre's formula its valuation is at least
\(j\mu-\log_p(M+1)\), since \(m+j\le bm\).
The valuation after multiplication by \(z^{r-j}\) is at least
\(re-\log_p(M+1)\).
If \(e>\mu\), the \(p\)-adic exponential converges and \(Q(z)\)
is the negative tail at a root of \(G\). In \eqref{eq:pade-beta}
put \(n=am+v\le(b-1)m\). Legendre's formula gives
\[
 v_p\!\left(\frac{n!}{(n+m+\ell)!}\right)
 \ge -(m+\ell)\mu-\mu s_p(n)
 \ge -(m+\ell)\mu-\log_p(M+1).
\]
After multiplication by \(z^{bm+\ell}\), the bound increases
with \(\ell\) and has its minimum at \(\ell=1\). In both cases,
\begin{equation}\label{eq:pade-local}
 v_p(Q(z))\ \ge\ (bm+1)e-N\min(e,\mu)-\log_p(M+1).
\end{equation}
For \(e=0\), integrality gives \(v_p(Q(z))\ge0\).

Now let \(z\) be a complex root and set \(R=|z|\).
Integrate each \(v!\) in \eqref{eq:pade-coeff} and use
the nonnegative coefficients of the polynomial in \(w\).
Also sum the beta tail before taking its absolute value.
These two operations give
\begin{equation}\label{eq:pade-integrals}
 \frac{|Q(z)|}{R^{bm+1}}\le\frac1{m!}\min\left\{
 \begin{aligned}
 &\int_0^\infty e^{-Ru}(1+u)^{am}((1+u)^d-1)^m\,du,\\
 &\int_0^1 e^{u\Re z}(1-u)^{am}(1-(1-u)^d)^m\,du
 \end{aligned}\right\}.
\end{equation}
In the second integral \(1-(1-u)^d\le du\). When \(am>\Re z\),
integrating \(u^m e^{-(am-\Re z)u}\) on the positive ray gives
the second gain in the lemma. When \(\Re z\le am\), the integral
is at most one. If \(R\le m\rho_0\), maximizing the exponent
\(u\rho_0+a\log(1-u)\) gives the upper bound
\(e^{m(\rho_0-a-a\log(\rho_0/a))}\). If \(R\ge m\rho_0\),
bound the first integrand using \(((1+u)^d-1)^m\le(1+u)^{dm}\)
and reserve \(e^{-u}\) for integration. The maximum of the
remaining exponent divided by \(m\) is
\(g_+(\rho_0-1/m)\), where \(g_+(v)=b\log(b/v)-b+v\).
Since \(\rho_0-b/2=a/2-s/14>19a/56>1/m\), the interval
\([\rho_0-1/m,\rho_0]\) lies in \((b/2,b)\), where
\(g_+'(v)=1-b/v>-1\). Therefore
\(g_+(\rho_0-1/m)<g_+(\rho_0)+1/m\) and
\(|Q(z)|<R^{bm+1}/L(z)^N\).
Stirling's lower bound gives the simpler gain once
\(0.01(m+1)+\tfrac12\log(2\pi/m)>0\).
This expression increases for \(m\ge199\), since its derivative
is \(0.01-1/(2m)\). At \(199\) it is positive by \(2\pi>6\)
and \(e^4>199/6\), using the first five terms of the exponential series.
When \(d=1\), use the first integral directly, with
\(((1+u)^d-1)^m=u^m\). The tangent inequality
\[
 \log(1+u)\le \log(1+v)-\frac{v}{1+v}+\frac{u}{1+v},
 \qquad v=a^{-1/2},
\]
bounds the constant exponent by \(m/2\), since
\(\log(1+v)-v/(1+v)\le v^2/2\). The coefficient of
\(-u\) is \(R-am/(1+v)>m\sqrt a/(1+v)\). Integrating
\(u^m\) against that exponential proves the additional gain.

Take the product of the complex bounds over the \(D\) roots
and compare with the prime factorization of \(\Res(G,Q)\) using
\eqref{eq:pade-local}. At a prime dividing \(K\) the sum of the
positive \(e\)'s is \(v_p(K)\), and the sum of \(\min(e,\mu)\)
is at most \(\min(v_p(K),D\mu)\). Cancellation of the common
term \((bm+1)\log K\) gives \eqref{eq:pade-norm}.
\end{proof}

For later estimates, the two losses defining \(g\), divided
by \(s\), are the functions
\begin{equation}\label{eq:pade-loss}
 K_+^2\int_0^1\frac{u\,du}{1-\alpha K_+u},
 \qquad
 K_-^2\int_0^1\frac{u\,du}{1-\alpha+\alpha K_-u}.
\end{equation}
The first integrand increases since \(K_+\) increases and
\(1-\alpha K_+u\) decreases. For the second, put
\(D_-=1-\alpha+\alpha K_-u\). Its derivative with respect to
\(\alpha\) is
\[
 \frac{uK_-}{D_-^2}
 \left(K_-(1-u/2)-\frac{1-\alpha}{7}\right),
\]
which is nonnegative for \(0\le u\le1\), since
\(K_-/2\ge25/112>1/7\) on \(0\le\alpha\le3/4\).
At \(\alpha=.03,.1,1/3,.5,.75\), upper bounds for both
functions are respectively \(.128,.134,.160,.176,7/32\).

We also need the following refinement of the local estimate.
Partial fractions in the tail formula yield
\[
 [z^{bm+\ell}]e^z B(z)=
 \frac{(-1)^m d^m}{(\ell-1)!}
 \sum_{j=1}^\ell
 \frac{(-1)^{j-1}\binom{\ell-1}{j-1}}
      {\prod_{i=0}^m(am+di+j)} .
\]
Indeed the coefficient of \(1/(x+j)\) in \(1/(x+1)_\ell\) is
\[
 \left(\prod_{\substack{1\le u\le\ell\\u\ne j}}(u-j)\right)^{-1}
 =\frac{(-1)^{j-1}}{(\ell-1)!}\binom{\ell-1}{j-1}.
\]
Apply this to \(x=am+di\) in the factorial expansion of the
tail. The remaining sum over \(i\) is
\(\sum_i(-1)^i\binom mi/(x+di)=d^m m!/\prod_{i=0}^m(x+di)\).
The coefficient of \(1/(x+di)\) on the right is
\(m!/\prod_{t\ne i}(t-i)=(-1)^i\binom mi\).
In both identities, matching the simple poles leaves a polynomial
that tends to zero at infinity.

For \(e>\mu\), the tail term indexed by \(\ell,j\) has
valuation at least
\[
 (bm+1)e+m v_p(d)
  -v_p\!\left(\prod_{i=0}^m(am+di+j)\right)+(\ell-1)(e-\mu),
\]
using \(v_p((\ell-1)!)\le(\ell-1)\mu\).
At fixed \(j\) the weakest bound is at \(\ell=j\).
Taking the infimum gives
\begin{equation}\label{eq:pade-progression}
 v_p(Q(z))\ge (bm+1)e+m v_p(d)-
 \sup_{j\ge1}\left\{v_p\!\left(\prod_{i=0}^m(am+di+j)\right)
                        -(j-1)(e-\mu)\right\}.
\end{equation}

\subsection{Rotations and the prime support}

Write \(U_i(j)\) for the reverse Newton ordinates
\[
 U_i(j)=v_p((bi-dj+1)_{dj})-v_p\binom ij,\qquad i=N-1,N.
\]
For \(L_p=\lfloor\log_p M\rfloor\), the \(dj\) factorial
arguments contain at most \(dj/p^\nu+1\) multiples of \(p^\nu\)
for \(1\le\nu\le L_p\), and none for \(\nu>L_p\).
Summing and dropping the binomial valuation gives
\(U_i(j)\le dj/(p-1)+L_p\). The concave hull has the
same affine upper bound, so its value at \(h\) bounds the
sum of the largest \(h\) valuations of roots of \(R_i\).
In particular
\begin{equation}\label{eq:norm-affine}
 v_p(F(0))\le H/(p-1)+L_p,\qquad
 p>2H\ \Longrightarrow\ v_p(F(0))\le L_p.
\end{equation}
Also every prime of either initial block \((ai+1)_H\)
divides \(F(0)\). Otherwise \(F\bmod p\), of degree \(h\) and
nonzero constant term, would divide a power of \(x\) times
a polynomial of degree less than \(h\).

\begin{lemma}[Cyclotomic restriction]\label{lem:kummer}
Above any one root of \(F\), the roots of \(G\) in
Lemma~\ref{lem:resultant} are a union of regular rotation cosets
of common size \(k_0\mid d\). If \(q=d/k_0\ge2\), and the
starting-point inequality of Lemma~\ref{lem:degree-cap} holds,
then
\begin{equation}\label{eq:kummer-support}
 2H<P=M^{1/(q+1)},\qquad p\mid F(0)\ \Longrightarrow\ p\le P.
\end{equation}
If \(k_0=d\), then \(k=d\).
\end{lemma}

\begin{proof}
Let \(x\) be a root of \(F\) and adjoin \(\mu_d\) to \(\Q(x)\).
Over this field the extension generated by \(z^d=-x\) is Galois,
and \(\sigma\mapsto\sigma(z)/z\) identifies the group with a
subgroup of \(\mu_d\), of order \(k_0\). Any other lift differs
from \(z\) by a root of unity already adjoined. This proves
the coset assertion and shows that \(z^{k_0}\) is a
\(q\)-th root of \(-x\) in \(\Q(x,\mu_d)\).
At \(p>d\) the cyclotomic extension is unramified, by lifting
roots of unity of order prime to \(p\) from the residue field.
Thus every positive integer-normalized local valuation of \(x\)
in \(\Q(x)\) is a multiple of \(q\). With normalization
\(v_p(p)=1\) it is at least \(q/h\), and \(v_p(F(0))\) is
also a multiple of \(q\).

Suppose \(p>H+d\), \(p^{q+1}>M\), and \(p\mid F(0)\).
For a selected slope \(\lambda\ge q/h\), the affine
upper bound puts every maximizing index of \(U_i(j)-\lambda j\)
below \(qh/(q-1)\le2h\). In this range \(j! \) is a unit.
The factorial jumps are spaced by at least \(h+1\) indices,
with heights at most \(q\), and the binomial valuation cannot
decrease. Thus the maximizing indices can only be zero and
the first factorial jump. By \eqref{eq:norm-affine}, the norm
exponent is exactly \(q\). The edge at each order must have
height \(q\), since it accounts for all selected positive slopes.
The equal widths then place the difference of the two reached
multiples of \(p^q\) in \([a+1,a+2d-1]\). Consequently every norm
prime is at most
\[
 \max\{H+d,\ M^{1/(q+1)},\ (a+2d-1)^{1/q}\}.
\]
The starting-point inequality supplies an initial-block prime
above \(2H\) by Lemma~\ref{lem:degree-cap}. If
\((a+2d-1)^{1/q}>2H\), then
\(s<(1+\alpha)/(4h^2)\), and \(b<dN\) gives
\[
 \frac{M}{(a+2d-1)^{3/2}}>
 \frac{1}{\sqrt s(1+\alpha)^{3/2}}>
 \frac{2h}{(1+\alpha)^2}>1.
\]
It follows that \(M^{1/(q+1)}>(a+2d-1)^{1/q}\) in this case.
If the last entry in the maximum is at most \(2H\), the
initial-block prime directly gives \(M^{1/(q+1)}>2H\).
Both cases prove \eqref{eq:kummer-support}.
\end{proof}

A coset of size \(n\ge2\) places at least \(\lfloor n/2\rfloor\)
points in the closed left half-plane, where the gain is at
least \(m(1-\alpha)/\alpha\). We will use a stronger average
when \(\alpha\) is small.

\begin{lemma}[Angular average]\label{lem:angles}
Assume \(N\ge200\) and \(\alpha<.03\), and put
\(g_0=1.01+.128s\) and \(C'=(1-\alpha)/(2\alpha)\).
In a coset of size \(n\ge2\), order the arguments
\(\theta_j\) by absolute value, and put \(u_j=|1-e^{i\theta_j}|^2\).
Let \(r\) be the number of \(j>0\) with \(u_j<4\alpha\).
The logarithm of the geometric mean of the gains is at least
\begin{equation}\label{eq:pade-angles}
 \log m+(1-1/n)\log C'
 +\frac{2\log(2n/\pi)-g_0}{n}
 -\frac{r}{n}(g_0+\log(2(1-\alpha)))-J,
\end{equation}
where
\[
 \frac rn\le \frac2\pi\arcsin\sqrt\alpha<.111,\qquad
 J=\frac{2\log2}{\pi}
        \bigl(\arcsin\sqrt\alpha+\sqrt{\alpha(1-\alpha)}\bigr)<.153.
\]
For \(n\le9\) one has \(r=0\) and may use
\(J=2(\log2)\alpha(n^2-1)/(3n)\). A weaker bound, used below,
takes \(C=C'/2\) in \eqref{eq:pade-angles}, replaces
\(\log(2(1-\alpha))\) by \(\log(1-\alpha)\), and omits \(J\).
\end{lemma}

\begin{proof}
Where \(u_j\ge4\alpha\), the gain is at least
\(mC'u_j(1-2\alpha/u_j)\). For nonnegative cosine use
\(|z|\le bm\) in \(am-\Re z\). For negative cosine use
\(|z|>am\). Use the baseline \(me^{-g_0}\) at the nearest
argument and at the other \(r\) arguments. The cyclotomic
product gives \(\prod_{j>0}u_j\ge(2n/\pi)^2\), because
\(|\theta_0|\le\pi/n\) and
\(|\sin(n\theta_0/2)/\sin(\theta_0/2)|\ge2n/\pi\).
The removed factors \(C'u_j\) are less than \(2(1-\alpha)\).
The angle spacing gives \(|\theta_j|\ge\pi j/n\) and the
bound on \(r/n\). For the retained factors, apply
\(\log(1-y)\ge-2(\log 2)y\) for \(0\le y\le1/2\).
The mean correction is bounded by the integral of the
decreasing function
\(2(\log2)\min\{1/2,\alpha/(2\sin^2(\pi x/2))\}\) over
\(0\le x\le1\), which equals \(J\).
If \(n\le9\), then \(\sin^2(\pi/(2n))>.03\) and the
discrete sum gives the claimed improvement using
\(\sum_{j=1}^{n-1}\csc^2(\pi j/(2n))=2(n^2-1)/3\).
Finally the difference between the strong and weak bounds
is \((1-(r+1)/n)\log2-J>0\).
\end{proof}

\subsection{Large primes in the norm}

\begin{lemma}[Two possible edge heights]\label{lem:two-heights}
Assume \(d\ge2\) and let \(p>2H+d\), \(p\mid F(0)\), and
\(p^4>M\). Suppose neither
interval \([bi-6H,bi]\), \(i=N-1,N\), contains a positive
multiple of \(p^3\). In particular this holds if \(p^3>M\).
The selected Newton edge at order \(i\) joins zero to the first
jump, of width \(w_i\) and height \(e_i\in\{1,2\}\).
Write \(u_i=bi-dw_i+q_i\), \(1\le q_i\le d\), for the reached
multiple of \(p\). The heights and widths satisfy
\[
 \frac{e_N}{w_N}=\frac{e_{N-1}}{w_{N-1}},\qquad
 v_p(F(0))\le\min(e_N,e_{N-1}).
\]
If the heights are equal, \(u_N-u_{N-1}\in[a+1,a+2d-1]\).
Otherwise \(v_p(F(0))=1\), the widths are in ratio two,
and \(p\) divides an integer in one of
\begin{equation}\label{eq:unequal-products}
 b(N-2)+[2-d,2d-1],\qquad b(N+1)+[2-d,2d-1].
\end{equation}
In this case \(p^2\) has a positive multiple in one of
\([bi-2H,bi]\). In both cases \(u_N\ne u_{N-1}\) and
\[
 u_N-u_{N-1}\in[b-H-d+1,b+H+d-1],
 \qquad p^{v_p(F(0))}\mid (u_N-u_{N-1}).
\]
\end{lemma}

\begin{proof}
The slope \(\lambda\) is at least \(1/h\). The affine
count with error three puts all maximizing indices below
\(6h\), and the absence of terminal cubes improves this to
\(4h<p\). The binomial valuation is nondecreasing there.
A factorial jump raises the ordinate by at most two, and
successive jumps are more than \(2h\) apart. Thus the
weighted ordinates \(U_i(j)-\lambda j\) strictly decrease
between jumps and from one jump index to the next. Only zero
and the first jump can maximize. This proves the first assertions.

The statements for equal heights follow by subtracting the
\(u_i\). For unequal heights the smaller width is the reduced
denominator of the slope and is at most \(h\). Twice the
reached multiple with the shorter width minus the other
belongs to \eqref{eq:unequal-products}. The wider edge
reaches a multiple of \(p^2\) within \(2H\) of the terminal.
In either case, \(|w_N-w_{N-1}|\le1/\lambda\le h\)
gives the stated interval for the difference. If the two
multiples were equal then \(w_N>w_{N-1}\), since \(b>d\).
The upper binomial numerator interval would contain the
lower one, with the same factorial valuation at the first
jump. This gives \(e_N\le e_{N-1}\), contradicting the common
positive slope. The norm bound on \(v_p(F(0))\) gives the last
divisibility.
\end{proof}

\begin{lemma}[Prime sums used below]\label{lem:prime-estimates}
Let \(S(x)=\sum_{p\le x}\log p/(p-1)\). Then for \(x>1\),
\[
 \pi(x)\le\frac{x}{\log x}\left(1+\frac{1.2762}{\log x}\right),
 \qquad S(x)<\log x+.2.
\]
For \(x\ge350\) the latter bound improves to \(\log x-.49\),
and for \(\log x\ge8\) to \(\log x-.54\).
\end{lemma}

\begin{proof}
The prime count is Dusart's Theorem 6.9 \cite{dusart-primes}.
Theorem 6.11 gives an upper bound for
\(\sum_{p\le x}(\log p)/p\) of
\(\log x+E+.2/\log x+.2/\log^2 x\) when \(x>2974\).
Add the prime-power correction using
\(E=-\gamma_{\rm E}-\sum_p\log p/(p(p-1))\) and
\(\gamma_{\rm E}>.57\). This proves the claimed bound
at \(\log x\ge8\). For \(1<x<3000\), the excess \(S(x)-\log x\)
decreases between primes. Enumerate primes by trial division
and bound the excess at each prime using the expansion
\(\log r=2\sum_{j=0}^9 z^{2j+1}/(2j+1)+\epsilon\), where
\(r\in[1,2]\), \(z=(r-1)/(r+1)\), and
\(0\le\epsilon\le2z^{21}/(21(1-z^2))\).
Normalize arguments by powers of two. This gives an upper
bound below \(.144\) on the excess, with its maximum at \(p=3\).
Starting at \(x=350\) and checking that endpoint as well gives
an upper bound below \(-.499\) through \(3000\).
\end{proof}

